%% PRÄAMBEL =====================================
\documentclass[
12pt,						% Schriftgröße
a4paper,					% Papierformat
oneside, 					% einseitiges (oneside) oder zweiseitiges (twoside) Dokument
listof=totoc, 				% Tabellen- und Abbildungsverzeichnis ins Inhaltsverzeichnis
bibliography=totoc,			% Literaturverzeichnis ins Inhaltsverzeichnis aufnehmen
titlepage, 					% Titlepage-Umgebung statt \maketitle
headsepline, 				% horizontale Linie unter Kolumnentitel
%abstracton,				% Überschrift beim Abstract einschalten, Abstract muss dazu in {abstract}-Umgebung stehen
DIV=12,						% Satzspiegeleinstellung, 12 ist Standard bei KOMA-Script
%BCOR=6mm,					% Bindekorrektur, die den Seitenspiegel um 6mm nach rechts verschiebt,
%cleardoublepage=empty,	% Stil einer leeren eingefügten Seite bei Kapitelwechsel
parskip,					% Absatzabstand bei Absatzwechsel einfügen
ngerman
]{scrbook}

%% PACKAGES UND EINSTELLUNGEN ===================

%% Debugging
%\usepackage{showframe}	% Anzeigen des Seitenlayouts; für Debugging auskommentieren
%\usepackage{lipsum}	% Erzeugen von Fülltext ('Lorem ipsum ...') zum ausprobieren von neuen Funktionen

\usepackage{scrhack}			% erlaubt Anpassung der Koma-Skript Klassen	
\usepackage[utf8]{inputenc} 	% ermöglicht die direkte Eingabe von Umlauten
\usepackage[T1]{fontenc} 		% Ausgabe aller zeichen in einer T1-Codierung (wichtig für die Ausgabe von Umlauten!)
\usepackage{babel} 				% deutsche Trennungsregeln und Übersetzung der festcodierten Überschriften (chapter => kapitel, usw.)

%% Einstellung der Abstände
\setlength{\parindent}{0ex} 	% bei neuem Abschnitt nicht einrücken
\usepackage[onehalfspacing]{setspace} % Zeilenabstand festlegen

% Abstände vor/nach neuen Abschnitten
\RedeclareSectionCommand[
  %runin=false,
  afterindent=false,
  beforeskip=0pt,
  afterskip=2\baselineskip]{chapter}
\RedeclareSectionCommand[
  %runin=false,
  afterindent=false,
  beforeskip=\baselineskip,
  afterskip=.5\baselineskip]{section}
\RedeclareSectionCommand[
  %runin=false,
  afterindent=false,
  beforeskip=.75\baselineskip,
  afterskip=.5\baselineskip]{subsection}

% -----------------------------------------------
% Folgende Einstellungen sind bei größeren Arbeiten mit viel Text zu empfehlen.
% Hierbei oben DIV=16 einstellen und Zeile \usepackage[onehalfspacing]{setspace} auskommentieren.
% \linespread{1.2}\selectfont     % Zeilenabstand erhöhen - größere Werte als 1.2 nicht verwenden!!
% -----------------------------------------------

%% Schriftarten
\usepackage{helvet}		% Änderung der Serifenlosen Schriftart auf Helvetica (Arial); kann bei Bedarf auskommentiert werden
\usepackage{mathptmx}	% Änderung der Serifen und Mathe Schriftart auf Times New Roman; kann auskommentiert werden

%% Zeichen und Mathematik
\usepackage{amsmath, amssymb}	% Beinhaltet Ergänzungen für Formeln
\usepackage{siunitx}		% Vereinfacht die Eingabe von Einheiten in Formeln
\sisetup{
	number-unit-product = \;,
	inter-unit-product = \:,
	exponent-product = \cdot,
	output-decimal-marker = {,},
	per-mode = fraction
}
\usepackage{textcomp} 			% Ermöglicht Einsatz von Sonderzeichen (z.B. Euro-Zeichen, ...)
\usepackage{eurosym}			% bessere Darstellung Euro-Symbol mit \euro

%% Einbindung von Grafiken
\usepackage{graphicx}  	% Einbinden von jpg, png und pdf Grafiken
\usepackage{float}			% Ermöglicht "Floating-Objects" im Dokument, also Objekte, die frei auf der Seite platziert werden können
\usepackage{pdfpages}			% Ermöglicht Einführung von bestehenden PDF Seiten (z.B. Datenblätter, Zeichnungen, ...)

% Bildunter- & überschriften
\usepackage[
format=hang,		% Formatierungseinstellungen von Unter- / Überschriften
font=normal,
labelfont=bf,
justification=RaggedRight,
singlelinecheck=true,
aboveskip=1mm
]{caption}

%% Einbinden von Literatureinträgen
\usepackage[
backend=biber, 		% Hilfsprogramm "biber" beim Compilieren nutzen (statt "biblatex" oder "bibtex")
style=alphabetic, 	% Zitierstil (siehe Dokumentation!!!)
natbib=true,		% Bereitstellen von natbib-kompatiblen Zitierkommandos
hyperref=true, 		% hyperref-Paket verwenden, um Links zu erstellen
]{biblatex}
\addbibresource{literature/literature.bib} 	% Einbinden der .bib-Datei

% Folgende Zeilen auskommentieren, falls runde Klammern und ein "vgl." bei Zitaten erscheinen sollen.
%\makeatletter
%\renewcommand{\@cite}[2]{(vgl. {#1\if@tempswa , #2\fi})} 
%\renewcommand{\@biblabel}[1]{(#1)}
%\makeatother

\usepackage{enumitem}			% Package für mehr Aufzählungsoptionen
\usepackage{tabto}				% Ermöglicht Verwendung von Tabstops zur Einrückung von Text

\usepackage[					% Einstellungen für Referenzierung
hyperfootnotes=false			% im pfd-Output werden Fußnoten nicht verlinkt
]{hyperref}

\usepackage[					% Einstellungen für Fußnoten
bottom,							% Ausrichtung  = unten
multiple,						% Trennung mehrerer Fußnoten durch einen Separator
hang,
marginal
]{footmisc}

%% weitere Verzeichnisse
%\usepackage[intoc]{nomencl} 	% zur Erstellung des Abkürzungsberzeichnisses (veraltet)
\usepackage[toc]{glossaries}
\makeglossaries

\usepackage{calc}			% Ermöglicht eine andere Eingabe von Längen (0.8\textwidth)
\usepackage{xcolor} 		% Ermöglicht Verwendung von Farben in nahezu allen Farbmodellen

% Benutzerdefinierte Farben erzeugen (werden später für Quellcode-Darstellung verwendet)
\definecolor{mygreen}{rgb}{0,0.6,0}
\definecolor{mygray}{rgb}{0.5,0.5,0.5}
\definecolor{mymauve}{rgb}{0.58,0,0.82}

% Zur Einbindung von Tikz und PGF-Plot-Grafiken
\usepackage{pgfplots}
\usepackage{pgfplotstable}
\pgfplotsset{compat=newest,width=0.6\linewidth}
\usepgfplotslibrary{smithchart} 	% Einbindung von Smith-Diagrammen (HF-Technik) möglich
\usepackage{tikz}					% Tikz sollte nach Listings Pakete geladen werden.
\usetikzlibrary{arrows}

%% Verbatim und Quellcode
\usepackage{verbatim}	% Eingabe von verbatim-Text mit den Umgebungen {verb} und \verb||
\usepackage{listings}			% Darstellung von Quellcode mit den Umgebungen {lstlisting}, \lstinline und \lstinputlisting
\lstset{literate=					% Ermöglicht Umlaute innerhalb von Listing-Umgebungen
	{Ö}{{\"O}}1
	{Ä}{{\"A}}1
	{Ü}{{\"U}}1
	{ß}{{\ss}}1
	{ü}{{\"u}}1
	{ä}{{\"a}}1
	{ö}{{\"o}}1
}

% Einstellung der Quellcode-Hervorhebung
\lstset{ %
	backgroundcolor=\color{white},   % choose the background color; you must add \usepackage{color} or \usepackage{xcolor}; should come as last argument
	basicstyle=\footnotesize,        % the size of the fonts that are used for the code
	breakatwhitespace=false,         % sets if automatic breaks should only happen at whitespace
	breaklines=true,                 % sets automatic line breaking
	captionpos=t,                    % sets the caption-position to (b) bottom or (t) top
	commentstyle=\color{mygreen},    % comment style
	deletekeywords={...},            % if you want to delete keywords from the given language
	escapeinside={\%*}{*)},          % if you want to add LaTeX within your code
	escapeinside={(*@}{@*)},
	extendedchars=true,              % lets you use non-ASCII characters; for 8-bits encodings only, does not work with UTF-8
	frame=none,	                   	% "single" adds a frame around the code; "none"
	keepspaces=true,                 % keeps spaces in text, useful for keeping indentation of code (possibly needs columns=flexible)
	keywordstyle=\color{blue},       % keyword style
	language=[LaTeX]TeX,             % the language of the code
	morekeywords={*,nomenclature},   % if you want to add more keywords to the set
	numbers=left,                    % where to put the line-numbers; possible values are (none, left, right)
	numbersep=5pt,                   % how far the line-numbers are from the code
	numberstyle=\tiny\color{mygray}, % the style that is used for the line-numbers
	rulecolor=\color{black},         % if not set, the frame-color may be changed on line-breaks within not-black text (e.g. comments (green here))
	showspaces=false,                % show spaces everywhere adding particular underscores; it overrides 'showstringspaces'
	showstringspaces=false,          % underline spaces within strings only
	showtabs=false,                  % show tabs within strings adding particular underscores
	stepnumber=1,                    % the step between two line-numbers. If it's 1, each line will be numbered
	stringstyle=\color{mymauve},     % string literal style
	tabsize=2,	                   % sets default tabsize to 2 spaces
	title=\lstname                   % show the filename of files included with \lstinputlisting; also try caption instead of title
}

% -----------------------------------------------
% Zum Aktualisieren des Abkürzungsverzeichnisses (Nomenklatur) bitte auf der Kommandozeile folgenden Befehl aufrufen :
% makeindex <Dateiname>.nlo -s nomencl.ist -o <Dateiname>.nls
% Oder besser: Kann in TexStudio unter Tools-Benutzer als Shortlink angelegt werden
% Konfiguration unter: Optionen-Erzeugen-Benutzerbefehle: makeindex -s nomencl.ist -t %.nlg -o %.nls %.nlo
% -----------------------------------------------

%% DEFINES ======================================

% Hier die persönlichen Daten eingeben:
% Die Befehle können als Variablen betrachtet werden, die im Titelblatt eingesetzt werden
% Die Felder können bei nichtverwendung leer gelassen werden
\newcommand{\titel}{Titel der wissenschaftlichen Arbeit}
\newcommand{\untertitel}{Untertitel (optional)}
\newcommand{\arbeit}{Projektbericht / Studienarbeit / Laborbericht / etc.}
\newcommand{\studiengang}{Studiengang}
\newcommand{\studienrichtung}{Studienrichtung}
\newcommand{\studienschwerpunkt}{Studienschwerpunkt}
\newcommand{\autor}{Name des Autors}
\newcommand{\matrikelnr}{Matrikelnummer}
\newcommand{\kurs}{Kursbezeichnung}
\newcommand{\firma}{Duales Partnerunternehmen}
\newcommand{\abgabe}{Datum der Abgabe}
\newcommand{\betreuerdhbw}{Betreuer der DHBW}
\newcommand{\betreuerfirma}{Betreuer des Partnerunternehmens}
\newcommand{\jahr}{20xx}			% für Angabe im Copyright-Vermerk der Titelseite

% Folgende Zeilen definieren Abkürzungen, um Befehle schneller eingeben zu können
\newcommand{\bs}{\textbackslash}
\newcommand*\diff{\mathop{}\!\mathrm{d}}				% Differentialzeichen
\newcommand*\Diff[1]{\mathop{}\!\mathrm{d^#1}} 		% Differentialzeichen höherer Ableitung
\newcommand*\jj{\mathop{}\!\mathrm{j}}					% Komplexe Zahl j (in E-Tech gebräuchlich)
\let\texteuro\euro % Eingabe des Euro-Zeichen über \texteuro oder \euro

% Abkürzungseinträge
% \newacronym{kürzel}{abkürzung}{langform}
\newacronym{zb}{z.B.}{zum Beispiel}
% Referenzierung
% \gls{kürzel}

%% BEGINN DES DOKUMENTENINHALTS =================

\begin{document}
\let\texteuro\euro					% Eingabe \texteuro, € oder \euro erzeugt gleiches Ergebnis
\setcounter{secnumdepth}{3}		% maximale Nummerierungstiefe fürs Inhaltsverzeichnis 
\setcounter{tocdepth}{3}
\sffamily							% für die Titelei wird serifenlose Schrift verwendet

%% TITELEI --------------------------------------

\pagenumbering{roman}				% kleine, römische Seitenzahlen für Titelei
\thispagestyle{plain}
\hypersetup{pageanchor=false}
\begin{titlepage}
\enlargethispage{20mm}
\sffamily 								% Serifenlose Grundschrift für die Titelseite einstellen

%% Firmenlogo
%\parbox{0.5\linewidth}{
\begin{flushleft}		% Text linksbündig ausrichten
	\includegraphics[width=0.25\textwidth]{images/Firmenlogo}\\[5ex] % Hier ggf. ein Logo der Firma einfügen
\end{flushleft}
%}
%
%% DHBW-Logo
%\parbox{0.5\linewidth}{
\begin{flushright}		% Text rechtsbündig ausrichten
	\includegraphics[width=0.25\textwidth]{images/DHBW_Logo}\\[5ex]
\end{flushright}
%}
%
%% Titel
\begin{center} 	% Text zentrieren
{\fontsize{20.74pt}{24pt}\selectfont
	\textbf{\titel}\\[1.5ex]}
%
%% Untertitel
{\fontsize{14pt}{17pt}\selectfont
	\textbf{\untertitel}\\[5ex]}
%
%% Arbeit
{\fontsize{17pt}{20pt}\selectfont
	\textbf{\arbeit}\\[2ex]}
%
%% Studiengang
{\fontsize{14pt}{17pt}\selectfont
	Studiengang \studiengang\\[2ex]}
%
%% Studienrichtung und Autor
{\fontsize{12pt}{14pt}\selectfont
	Studienrichtung \studienrichtung\\[1ex]
	Duale Hochschule Baden-Württemberg Ravensburg, Campus Friedrichshafen\\[5ex]
	von\\[1ex]
	\autor\\[15ex]}
\end{center}
%
%% weitere Angaben
\begin{flushleft}
{\fontsize{12pt}{14pt}\selectfont
\begin{tabular}{ll}		% 
	Abgabedatum:			& \quad \abgabe \\
	%Bearbeitungszeitraum:	& \quad 12 Wochen   \\
	Matrikelnummer: 		& \quad \matrikelnr \\
	Kurs: 					& \quad \kurs \\
	%Ausbildungsfirma:	 	& \quad \firma \\
	Abgabe bei: 			& \quad \betreuerdhbw \\ [2ex]
\end{tabular}
}
\end{flushleft}
%
%%%%% Nachfolgende Zeilen einkommentieren, wenn Copyrightvermerk gewünscht ist
%\begin{flushleft}
%{\fontsize{11pt}{13pt}\selectfont
%Copyrightvermerk:\\
%Dieses Werk einschließlich seiner Teile ist \textbf{urheberrechtlich geschützt}. Jede Verwertung außerhalb der engen Grenzen des Urheberrechtgesetzes ist ohne Zustimmung des Autors unzulässig und strafbar. Das gilt insbesondere für Vervielfältigungen, Übersetzungen, Mikroverfilmungen sowie die Einspeicherung und Verarbeitung in elektronischen Systemen.
%}
%\end{flushleft}
%\begin{flushright}
%{\fontsize{11pt}{13pt}\selectfont \copyright{} \jahr }
%\end{flushright}
%
\end{titlepage}
%
\cleardoublepage
\hypersetup{pageanchor=true}
 		% erzeugt die Titelseite
%% Ggf. folgende Zeile auskommentieren, falls der Sperrvermerk gewünscht ist.
\chapter*{Sperrvermerk} %*-Variante sorgt dafür, das der Sperrvermerk nicht im Inhaltsverzeichnis auftaucht
gemäß Ziffer 1.1.13 der Anlage 1 zu §§ 3, 4 und 5  der Studien- und Prüfungsordnung für die Bachelorstudiengänge im Studienbereich Technik der Dualen Hochschule Baden-Württemberg vom 29.09.2017 in der Fassung vom 25.07.2018:

Der Inhalt dieser Arbeit darf weder als Ganzes noch in Auszügen Personen außerhalb des Prüfungsprozesses und des Evaluationsverfahrens zugänglich gemacht werden, sofern keine anders lautende Genehmigung vom Dualen Partner vorliegt.\\[6ex]

Ort, den xx.xx.20xx \\[1ex]

\rule[-0.2cm]{5cm}{0.5pt} \\

\autor
\clearpage

\chapter*{Erklärung} %*-Variante sorgt dafür, dass die Erklärung nicht im Inhaltsverzeichnis auftaucht

gemäß Ziffer 1.1.13 der Anlage 1 zu §§ 3, 4 und 5  der Studien- und Prüfungsordnung für die Bachelorstudiengänge im Studienbereich Technik der Dualen Hochschule Baden-Württemberg vom 29.09.2017 in der Fassung vom 25.07.2018.

Ich versichere hiermit, dass ich meine Projektarbeit mit dem Thema:

\begin{quote}
	\textit{\titel}
\end{quote}

selbstständig verfasst und keine anderen als die angegebenen Quellen und Hilfsmittel benutzt habe. Ich versichere zudem, dass die eingereichte elektronische Fassung mit der gedruckten Fassung übereinstimmt.\\[6ex]

Ort, den xx.xx.20xx \\[1ex]

\rule[-0.2cm]{5cm}{0.5pt} \\

\autor
\clearpage

\rmfamily

%\thispagestyle{empty}
%\cleardoublepage 		% Einbinden der eidestattlichen Erklärung und Sperrvermerk
\chapter*{Kurzfassung} %*-Variante sorgt dafür, das Abstract nicht im Inhaltsverzeichnis auftaucht
\label{cha:abstract_ger}
Im Rahmen der ersten und zweiten Praxisphase wurde ein erster Einblick in das Geschäftsgeschehen möglich. So ergab es sich, die Tätigkeit eines Projektmanagers und Planungsingenieurs im Ansatz kennenzulernen. 

Ein Großteil der ersten Praxisphase wurde von mir damit verbracht, ein Excel-Tool zur automatischen Erstellung eines Planverzeichnisses nach betriebsinternen Vorgaben zu entwickeln, womit interne Geschäftsprozesse vereinfacht und beschleunigt werden sollen. Dieses Tool sollte Teil einer weitergehenden Software-Umgebung zur Vereinfachung, Beschleunigung und Vereinheitlichung interner Geschäftsprozesse werden.

Auf Grundlage der dadurch angeeigneten Erfahrungen und Fähigkeiten in der Entwicklung von Excel-Arbeitshilfen bot sich die Möglichkeit, ein Planungstool für Beschallungsanlagen zu entwickeln. Das Tool soll den Fachplaner bei der Berechnung von Leitungsquerschnitten einzelner Lautsprecherkreise der 100-Volt-Technik unterstützen. Ziel ist es, den Planungsprozess zu vereinfachen und zu beschleunigen.

\chapter*{Abstract}
\label{cha:abstract_eng}
The first and second practice semester gave a first impression on the business activities and duties. Because of the right circumstances, it was possible for me to get to know the basic activities of a project manager and a planning engineer.

A considerable part of my first semester was spent by me developing an Excel tool for automatically creating a plan directory with internal corporate specifications, the purpose of this being to simplify and accelerate internal business processes. Additionaly, this tool was planned to be part of a wider software tool environment for simplification, acceleration and unifying purposes of planning processes.

Based on the collected experience and skills by developing such a tool, the opportunity to create a planning tool for public adress systems arose. The tool is designed to help the specialised planning engineer with calculating wire line cross-sections of single speaker circuits in the subject of 100-volt-technology. The goal is to simplify and accelerate the planning process.

\cleardoublepage
   		% Einbinden des Abstracts (der Kurzfassung)

\tableofcontents					% Erzeugen des Inhaltsverzeichnisses
\cleardoublepage

%% BEGINN DES DOKUMENTS -------------------------

\pagenumbering{arabic}					% arabische Seitenzahlen für den Hauptteil
\rmfamily

\chapter{Einleitung}
\label{cha:einleitung}
Problemstellung, Ziel und Vorgehensweise der Arbeit (Gegenstand und Ziele der Arbeit/Aufgabenbeschreibung, geplante Vorgehensweise, Einführung in Thema, Motivation der Aufgabenstellung/Voraus-
blick)

Die Einleitung soll den Ausgangspunkt der Arbeit umreißen, in kurzer Form zur Problemstellung hinführen und das Interesse der lesenden Person für die Arbeit wecken. Allgemeine Einleitung ins Thema, keine Unternehmens- oder Produktbeschreibungen, Organigramme u.ä., wenn diese nicht direkt zum Thema führen. Ziele
und Vorgehensweise nicht vermischen.
\chapter{Grundlagen}
\label{cha:grundlagen}
Ausgehend von der Aufgabenstellung ist der derzeitige Stand der Technik für die Lösungsfindung zu beschreiben. Es sind z.B. die Vor- und Nachteile bisheriger Lösungen bzw. fundamentaler Lösungsprinzipien fundiert
von und ggf. anderen Quellen darzulegen.


\chapter{Hauptteil}
\label{cha:hauptteil}
Anforderungsdefinition, Anforderungsanalyse, Lösungsgenerierung, Lösungsbewertung, Umsetzung \\
ggf. in mehreren sinnvollen Gliederungspunkten

Der Text soll knapp und klar sein und die wesentlichen Gedanken der Arbeit beinhalten. Ein gewähltes Verfahren oder ein bestimmter Lösungsweg muss begründet werden. Es ist nicht notwendig, alle Vorversuche einzeln zu schildern. Bei Versuchen sind Voraussetzungen und Vernachlässigungen sowie die Anordnung, Leistungsfähigkeit und Messgenauigkeit der Versuchsanordnung anzugeben.

Die Ergebnisse der Arbeit sind unter Berücksichtigung der Voraussetzungen ausführlich zu diskutieren und
mit den bereits bekannten Anschauungen und Erfahrungen zu vergleichen.

Ziel der Arbeit ist es, eindeutige Folgerungen und Richtlinien für die Praxis zu finden.

\begin{itemize}
	\item Anforderungsdefinition
	\item Anforderungsanalyse
	\item Lösungsgenerierung
	\item Lösungsbewertung
	\item Umsetzung
\end{itemize}
\chapter{Kritische Reflexion und Ausblick}
\label{cha:reflexion}
Aufgabenstellung, Vorgehensweise und wesentliche Ergebnisse werden kurz und präzise dargestellt und kritisch reflektiert. Die Zusammenfassung ist eigenständig verständlich. Länge ca. 1 bis 1,5 Seiten (Problem, Ziele, Vorgehensweise, Ergebnisse und Ausblick).

%% LITERATURVERZEICHNIS -------------------------

\interlinepenalty 10000				% Verhindert einen Umbruch mitten in Literatureinträgen
\nocite{*}								% Literaturverzeichnis ohne Zitation einblenden
\printbibliography						% Erstellen des Literaturverzeichnisses

%% VERZEICHNISSE --------------------------------

\setlength{\parskip}{0.5\baselineskip}

%% alle Abkürzungen, die in der Arbeit verwendet werden. Die Alphabetische Sortierung übernimmt Latex. Nachfolgend sind Beispiele genannt, welche nach Bedarf angepasst, gelöscht oder ergänzt werden können.

% Bei den unten stehenden Formelzeichen ist erläutert, wie explizite Sortierschlüssel über den Inhalt der eckigen Klammer angegeben werden.

% Allgemeine Abkürzungen %%%%%%%%%%%%%%%%%%%%%%%%%%%%
\nomenclature{DHBW}{Duale Hochschule Baden-Württemberg}
\nomenclature{bzw.}{beziehungsweise}
\nomenclature{etc.}{et cetera}
\nomenclature{z. B.}{zum Beispiel}
\nomenclature{o.Ä}{oder Ähnliches}

\nomenclature{PA}{Public Adress}
\nomenclature{ELA}{Elektroakustische Anlagen}
\nomenclature{PA}{Public Adress}
\nomenclature{SAA}{Sprachalarmanlagen}
\nomenclature{dB}{Dezibel}
\nomenclature{SPL}{Sound Pressure Level}
\nomenclature{THD}{Total Harmonic Distortion}
\nomenclature{HE}{Höheneinheit}
\nomenclature{DB KT}{Deutsche Bahn Kommunikationstechnik GmbH}
\nomenclature{TK}{Telekommunikation}

% Formelzeichen %%%%%%%%%%%%%%%%%%%%%%%%%%%%%%%%%%%%

\nomenclature[P]{$P$}{Leistung}
\nomenclature[U]{$U$}{Spannung}
\nomenclature[I]{$I$}{Stromstärke}
\nomenclature[R]{$R$}{Widerstand}

\nomenclature[Ls]{$L_s$}{Schallpegel}
\nomenclature[r]{$r$}{Reflektionskoeffizient}
\nomenclature[alpha]{$\alpha$}{Absorptionskoeffizient}
\nomenclature[T]{$T$}{Nachhallzeit eines Raumes}

\nomenclature[V]{$V$}{Gesamtvolumen des Raums}
\nomenclature[A]{$A$}{Gesamtabsorption}
\nomenclature[A]{$A_p$}{Absorption pro Person}
\nomenclature[S]{$S$}{Gesamtoberfläche}
\nomenclature[Si]{$S_i$}{Oberfläche}
\nomenclature[alphai]{$\alpha_i$}{Gesamtabsorption}
\nomenclature[m]{$m$}{Luftdruck bzw. Dämpfungskonstante}
\nomenclature[n]{$n$}{Anzahl der Personen im Raum}

\nomenclature[rho]{$\rho_{Cu}$}{Spezifischer Widerstand von Kupfer}
\nomenclature[L]{$L$}{Leiterlänge}
\nomenclature[A_min]{$A_{min}$}{Mindestquerschnittsfläche}
\nomenclature[d_min]{$d_{min}$}{Mindestdurchmesser}

% Dateiendungen %%%%%%%%%%%%%%%%%%%%%%%%%%%%%%%%%%%%
%\nomenclature{EMF}{Enhanced Metafile}
%\nomenclature{JPG}{Joint Photographic Experts Group}
%\nomenclature{PDF}{Portable Document Format}
%\nomenclature{PNG}{Portable Network Graphics}
%\nomenclature{XML}{Extensible Markup Language}




				% Datei mit Abkürzungen laden
%\renewcommand{\nomname}{Verzeichnis verwendeter Formelzeichen und Abkürzungen}
%\setlength{\nomlabelwidth}{.20\hsize}
%\renewcommand{\nomlabel}[1]{#1 \dotfill}
%\setlength{\nomitemsep}{-\parsep}
%\printnomenclature						% Erzeugen des Abkürzungsverzeichnises, siehe auch Inhalt der Datei pages/abkuerzungen.tex
\cleardoublepage

%\renewcommand{\glossaryname}{Glossar} % Glossar erzeugen

\printglossaries
\cleardoublepage

\listoffigures 							% Erzeugen des Abbildungsverzeichnisses 
\cleardoublepage

\listoftables 							% Erzeugen des Tabellenverzeichnisses
\cleardoublepage

%% ANHANG ---------------------------------------

\appendix
%\clearpage
%\pagenumbering{Roman}					% große, römische Seitenzahlen für Anhang, falls gewünscht
\chapter{Allgemeine Ergänzungen}
Hier sollen, wenn dies notwendig erscheint, Vorversuche, mathematische Herleitungen, Arbeitsanleitungen, Rechnerprogramme usw. ausführlich wiedergegeben werden. Ebenso können umfangreiche Ergebnisse gleichartiger Versuche in Diagramm- oder Tabellenform zusammengefasst werden.
Die im Anhang aufgenommenen Arbeiten sind in einem speziellen Anhang-Verzeichnis aufzulisten.

\section{Details zu bestimmten theoretischen Grundlagen}

\chapter{Technische Zeichnungen und Schaltpläne}

\clearpage


\end{document}
